\begin{table}[!t]
\caption{Sensitivity of the Risk Component's Association with Announcement Returns}
\label{tab:risk}
\centering
\footnotesize
\setlength{\tabcolsep}{3pt}
\begin{tabular}{lcc}
\toprule
\multicolumn{3}{l}{\textit{A. Benchmark and window: AUC (two-sided $p$)}} \\
\textbf{Outcome} & \textbf{Primary 55} & \textbf{Other 35} \\
\midrule
SPY, reaction day (primary) & 0.699 (0.04) & 0.724 (0.04) \\
Bank-sector index, reaction day & 0.514 (0.90) & 0.561 (0.66) \\
Market model, reaction day & 0.644 (0.14) & 0.637 (0.22) \\
SPY, $[0,+1]$ around release date & 0.423 (0.47) & 0.572 (0.57) \\
SPY, $[-1,+1]$ around release date & 0.343 (0.13) & 0.562 (0.62) \\
\midrule
\multicolumn{3}{l}{\textit{B. Influence in the primary panel: AUC}} \\
\textbf{Sample} & \textbf{SPY} & \textbf{Bank index} \\
\midrule
All 55 transactions & 0.699 & 0.514 \\
Leave one event out (range) & 0.639--0.731 & 0.494--0.573 \\
Without SVB, First Republic, WaMu & 0.736 & 0.540 \\
Without the five largest $|CAR|$ events & 0.733 & 0.518 \\
Announced 2008--2010 ($n=37$) & 0.717 & 0.442 \\
Announced 2011--2024 ($n=18$) & 0.679 & 0.711 \\
\bottomrule
\end{tabular}
\par\vspace{2pt}\parbox{0.94\columnwidth}{\scriptsize AUC of the Risk score $R$ against a positive CAR; $p$ from 10,000 permutations of
outcomes across event clusters. Normalization bounds and thresholds are fitted on the primary panel. The bank-sector index is the
equal-weighted return of 27 U.S.-listed bank holding companies, excluding the acquirer and banks with their own announcement
within four days. Panel B drops, in turn, each of the 51 event clusters, the three largest transactions, and the five event
clusters with the largest absolute CAR under each benchmark, and splits the panel by announcement year.}
\end{table}
